% Style follows ../class2/class2PT26.tex.
% Build in this directory (run twice to resolve page references and bookmarks):
% pdflatex -interaction=nonstopmode -halt-on-error class8PT26.tex
% Topics follow lectureNote.tex, sections 8 and 9; Levy's theorem is recalled from section 7.
\documentclass[a4paper,10pt,pdflatex,ja=standard,jafont=ipaex-type1,
  magstyle=real,scale=1]{bxjsarticle}
\usepackage{lmodern}
\usepackage{amsmath,amssymb,mathtools}
\usepackage{geometry}
\geometry{reset,left=22mm,right=22mm,top=20mm,bottom=17mm,
  headheight=13pt,headsep=6mm,footskip=10mm}
\usepackage{xcolor}
\usepackage{enumitem}
\usepackage{array,booktabs,tabularx}
\usepackage{fancyhdr}
\usepackage{lastpage}
\usepackage{hyperref}

\definecolor{ink}{HTML}{193A5C}
\definecolor{accent}{HTML}{237D82}
\definecolor{soft}{HTML}{EAF3F3}
\definecolor{muted}{HTML}{52616B}
\newcommand{\coursename}{確率論と統計学 I}
\newcommand{\lecturenumber}{8}
\hypersetup{unicode=true,colorlinks=true,urlcolor=accent,linkcolor=ink,
  pdfauthor={中島秀太},pdftitle={第8・9回：確率変数列の収束、分布の弱収束と極限定理}}
\pdfstringdefDisableCommands{\def\quad{ }}
\pagestyle{fancy}
\fancyhf{}
\fancyhead[L]{\small\color{ink}\coursename{}・授業要約と演習}
\fancyhead[R]{\small\color{ink}第\lecturenumber 回}
\fancyfoot[C]{\small\color{ink}\thepage/\pageref*{LastPage}}
\renewcommand{\headrulewidth}{0.3pt}
\setlength{\parindent}{0pt}
\setlength{\parskip}{2pt}
\setlength{\abovedisplayskip}{4pt}
\setlength{\belowdisplayskip}{4pt}
\setlength{\abovedisplayshortskip}{3pt}
\setlength{\belowdisplayshortskip}{4pt}
\setlist[itemize]{leftmargin=1.8em,itemsep=1pt,topsep=2pt,parsep=0pt}
\setlist[enumerate]{leftmargin=2em,itemsep=3pt,topsep=3pt,parsep=0pt}
\newcommand{\topic}[1]{\par\vspace{5pt}
  {\large\sffamily\mdseries\color{accent}#1}\par\vspace{2pt}}
\newcommand{\handouttitle}[1]{%
  \pdfbookmark[0]{#1}{lecture\lecturenumber}%
  \begin{center}
    {\Large\color{ink}#1}\par\vspace{3mm}
    {\color{accent}\rule{0.88\linewidth}{0.7pt}}
  \end{center}\vspace{-2mm}}
\newcommand{\keybox}[1]{\par\vspace{3pt}\begingroup
  \setlength{\fboxsep}{5pt}
  \setlength{\fboxrule}{0pt}
  \fbox{\parbox{\dimexpr\linewidth-2\fboxsep\relax}{#1}}%
  \endgroup\par\vspace{2pt}}
\newcommand{\examplelabel}[1]{{\sffamily\mdseries\color{ink}#1}}
\newcommand{\F}{\mathcal F}
\newcommand{\R}{\mathbb R}
\newcommand{\E}{\mathbb E}
\newcommand{\ind}{\mathrm{1}}

\begin{document}
\fontsize{10pt}{13.5pt}\selectfont
\handouttitle{第8回\quad 確率変数列の収束}

\topic{概収束と確率収束}
$X_n,X$ を同じ確率空間 $(\Omega,\F,P)$ 上の実数値確率変数とする。
\keybox{
$$\begin{aligned}
 X_n\to X\ \text{a.s.}
 &\quad\overset{\rm def}{\Leftrightarrow}\quad P(\{\omega:X_n(\omega)\to X(\omega)\})=1,\\
 X_n\xrightarrow{P}X
 &\quad\overset{\rm def}{\Leftrightarrow}\quad P(|X_n-X|>\varepsilon)\to0
 \quad(\forall\varepsilon>0).
\end{aligned}$$}
概収束では一つの標本を固定して数列を見る。確率収束では各時点の誤差確率を見る。

\examplelabel{定理}\quad $X_n$ が $X$ に概収束するならば、$X_n$ は $X$ に確率収束する。

\examplelabel{証明}\quad $\varepsilon>0$ を任意に固定し、$A_n=\{|X_n-X|>\varepsilon\}$ とおく。
概収束の仮定より、ほとんどすべての $\omega$ で $X_n(\omega)\to X(\omega)$ となる。
そのような $\omega$ では、十分大きな $n$ に対して $|X_n(\omega)-X(\omega)|\le\varepsilon$ となるので、
$\ind_{A_n}(\omega)\to0$ である。$0\le\ind_{A_n}\le1$ で定数関数 $1$ は可積分だから、
優収束定理より $P(A_n)=\E[\ind_{A_n}]\to0$ を得る。よって $X_n\xrightarrow{P}X$ である。

\topic{逆向きの反例：動く区間}
$\Omega=[0,1)$ に一様確率を入れ、$n=2^m+j$（$m\ge0$、$0\le j<2^m$）に対して
$$X_{2^m+j}(\omega)=\ind_{[j2^{-m},(j+1)2^{-m})}(\omega)$$
と定める。$0<\varepsilon<1$ なら $P(|X_{2^m+j}|>\varepsilon)=2^{-m}\to0$ なので
$X_n\xrightarrow{P}0$。しかし、どの $\omega$ でも各段で $1$ を一度とり、
$0$ も無限回とるため、列全体はどの標本でも収束しない。

\topic{Borel--Cantelli の補題}
事象列 $(A_n)_{n\ge1}$ について、次が成り立つ。
$$P(A_n\text{ が無限回起こる})=
 \begin{cases}
  0,&\sum_nP(A_n)<\infty,\\
  1,&\sum_nP(A_n)=\infty\text{ かつ }(A_n)\text{ は相互独立}.
 \end{cases}$$

\topic{概収束する部分列}
\examplelabel{定理}\quad $X_n$ が $X$ に確率収束するならば、$X$ に概収束する部分列 $(X_{n_k})$ が存在する。

\examplelabel{証明}\quad 各 $k\ge1$ について、確率収束より $P(|X_n-X|>2^{-k})\to0$（$n\to\infty$）である。
したがって $P(|X_{n_k}-X|>2^{-k})\le2^{-k}$ を満たす添字を $n_1<n_2<\cdots$ と順に選べる。
$A_k=\{|X_{n_k}-X|>2^{-k}\}$ とおくと、
$$\sum_{k=1}^{\infty}P(A_k)\le\sum_{k=1}^{\infty}2^{-k}=1.$$
Borel--Cantelli の補題より $P(A_k\text{ が無限回起こる})=0$ である。
よって、ほとんどすべての $\omega$ で、十分大きな $k$ に対して
$|X_{n_k}(\omega)-X(\omega)|\le2^{-k}$ となる。$2^{-k}\to0$ だから、$X_{n_k}$ は $X$ に概収束する。

\topic{演習問題}
\begin{enumerate}[label=\arabic*.]
\item $U$ は区間 $(0,1)$ 上の一様分布に従うとし、$X_n=U^n$、$Y_n=nU^n$ とおく。
それぞれの概収束と確率収束を調べ、期待値の極限も計算せよ。
\item $X_n$ が $X$ に確率収束し、確率 $1$ で $X_1\le X_2\le\cdots$ が成り立つとする。
$X$ に概収束する部分列を用いて、列全体が $X$ に概収束することを証明せよ。
\end{enumerate}

\topic{レポート問題}
\begin{enumerate}[label=\arabic*.]
\item 同じ確率空間上で $X_n\xrightarrow{P}X$、$Y_n\xrightarrow{P}Y$ とする。
確率収束の定義を用いて、$X_n+Y_n\xrightarrow{P}X+Y$ を証明せよ。
\item 独立なベルヌーイ確率変数 $B_n$ が $P(B_n=1)=1/(n+1)$ を満たすとする。
$B_n\xrightarrow{P}0$ を示せ。
独立な事象に対する Borel--Cantelli の補題を用い、$B_n$ が概収束しないことを示せ。
\end{enumerate}

\clearpage
\renewcommand{\lecturenumber}{9}
\handouttitle{第9回\quad 分布の弱収束と極限定理}

\topic{弱収束と分布収束}
$\mu_n,\mu$ を $\R$ 上のボレル確率測度、$C_b(\R)$ を有界連続関数全体とする。
$$\mu_n\Rightarrow\mu\quad\overset{de}{\Leftrightarrow}\quad
 \int_\R f\,d\mu_n\to\int_\R f\,d\mu\quad(\forall f\in C_b(\R)).$$
$X_n\Rightarrow X$ は分布 $\mu_{X_n}\Rightarrow\mu_X$ を意味する。
変数が同じ確率空間にあることは必要ない。同じ空間上では
$$X_n\to X\ \text{a.s.}\quad\Longrightarrow\quad
 X_n\xrightarrow{P}X\quad\Longrightarrow\quad X_n\Rightarrow X.$$
弱収束は、次の開集合の条件、閉集合の条件のそれぞれと同値である。
$$\mu(O)\le\liminf_n\mu_n(O)\quad(O\text{ は開}),\qquad
 \limsup_n\mu_n(F)\le\mu(F)\quad(F\text{ は閉}).$$

\topic{L\'evy の連続性定理}
$\varphi_n(t)=\int_\R e^{itx}\,\mu_n(dx)$、
$\varphi_\mu(t)=\int_\R e^{itx}\,\mu(dx)$ とする。
\keybox{
$$\mu_n\Rightarrow\mu\quad\Longleftrightarrow\quad
 \varphi_n(t)\to\varphi_\mu(t)\quad(\forall t\in\R).$$
さらに $\varphi_n(t)\to\varphi(t)$ がすべての $t\in\R$ で成り立ち、
$\varphi$ が原点で連続なら、$\varphi=\varphi_\mu$ となる確率測度 $\mu$ が一意に存在し、
$\mu_n\Rightarrow\mu$ となる。}
原点での連続性の必要性は次からわかる。$X_n\sim N(0,n)$ では
$$\varphi_n(t)=e^{-nt^2/2}\longrightarrow
 \begin{cases}1,&t=0,\\0,&t\ne0,\end{cases}$$
となり、極限は特性関数にならない。

\topic{緊密性と弱収束する部分列}
実数値確率変数列 $(X_n)$ が緊密であるとは
$$\lim_{R\to\infty}\sup_n P(|X_n|>R)=0$$
となることである。各確率は $X_n$ が定義された確率空間でとる。
緊密な確率変数列の任意の部分列には、さらに分布収束する部分列がある（Prokhorov の定理）。
すべての分布収束する部分列の極限分布が同じ $\mu$ なら、列全体の分布が $\mu$ に弱収束する。

\topic{演習問題}
\begin{enumerate}[label=\arabic*.]
\item $\sup_n\E[X_n^2]<\infty$ なら確率変数列 $(X_n)$ が緊密であることを示せ。
\item 各 $n\ge1$ に対し、$X_n$ は区間 $(-n,n)$ 上の一様分布に従うとする。
確率変数列 $(X_n)$ が緊密でないことを示せ。
\end{enumerate}

\topic{レポート問題}
\begin{enumerate}[label=\arabic*.]
\item $X_n,Y_n$（$n\ge1$）および $X,Y$ は、すべて同じ確率空間上の実数値確率変数とする。
$X_n\Rightarrow X$ と $Y_n\Rightarrow Y$ だけでは、$X_n+Y_n\Rightarrow X+Y$ とは限らないことを示す反例を構成せよ。
\item 緊密な確率変数列 $(Z_n)$ に対し、$Z_n$ の分布関数を $F_n$ とする。
部分列 $(F_{n_k})$ がすべての $q\in\mathbb Q$ で $F_{n_k}(q)\to G(q)$ を満たすと仮定し、
$x\in\R$ に対して $F(x)=\inf\{G(q):q\in\mathbb Q,\ q>x\}$ とおく。
\begin{enumerate}[label=(\arabic*),leftmargin=2.1em]
\item $F$ は単調非減少かつ右連続であり、$F$ の任意の連続点 $x$ で $F_{n_k}(x)\to F(x)$ が成り立つ。
\item 緊密性より $\lim_{x\to-\infty}F(x)=0$、$\lim_{x\to\infty}F(x)=1$ を示し、$F$ が分布関数であることを導け。
\item 一様確率を入れた $(0,1)$ 上で、$X_n(u)=\inf\{x\in\R:F_n(x)\ge u\}$、
$X(u)=\inf\{x\in\R:F(x)\ge u\}$ とおく。$X_n$ と $X$ の分布関数がそれぞれ $F_n$ と $F$ であり、
$X_{n_k}$ が $X$ に概収束することを示せ。さらに各 $f\in C_b(\R)$ について
$\E[f(X_{n_k})]\to\E[f(X)]$ を証明せよ。
\end{enumerate}
\end{enumerate}

\end{document}
